\chapter{Interference and Limitations --- Why Your Test Might Be Wrong}
\label{app:interference}

\clearpage
\section{Introduction}

Here is an uncomfortable truth about lab testing: many things can make your results wrong. Not slightly off --- completely wrong. Your medications can fake an abnormal result. That morning jog before your blood draw can change your potassium level. The supplement you take every day can make your thyroid look abnormal when it is perfectly fine.

This appendix is your guide to the things that can fool your lab tests. Read this before you panic about an abnormal result.

\clearpage
\section{Foods That Mess Up Your Tests}

What you eat --- and when you eat it --- has a direct effect on your lab results.

\subsection{What Preparation Is Needed}

\textbf{Tests that require fasting:}
\begin{itemize}
    \item \textbf{Fasting glucose:} Eating raises your blood sugar. A non-fasting glucose result is meaningless for diabetes screening.
    \item \textbf{Lipid panel:} Triglycerides can increase by 20--30\% after a meal. If you eat before a lipid panel, your triglycerides will be falsely elevated and your HDL may be falsely low.
    \item \textbf{Iron studies:} Serum iron drops after meals and can give misleading results for iron deficiency evaluation.
    \item \textbf{TIBC (total iron-binding capacity):} Affected by iron intake.
\end{itemize}

\textbf{How long to fast:} Most tests require 8--12 hours of fasting. Water is fine. Black coffee (without cream or sugar) is usually acceptable. Nothing else.

\subsection{Contraindications and Precautions}

\begin{longtable}{@{}p{4.5cm}p{10.5cm}@{}}
\toprule
\textbf{Food/Substance} & \textbf{Tests Affected} \\
\midrule
High-fat meal & Lipids (triglycerides falsely elevated), lipemic serum interferes with many tests \\
Bananas, avocados, oranges & Potassium (can transiently raise potassium) \\
Red meat & Uric acid (can elevate uric acid levels) \\
Coffee (caffeinated) & Cortisol, gastrin, catecholamines, free fatty acids \\
Alcohol & Liver enzymes (AST, ALT, GGT), triglycerides, uric acid, lactate \\
Spinach, beets & occult blood in stool (false positive) \\
Vitamin C (large doses) & Glucose (glucometer falsely low), creatinine, occult blood \\
Carrots, oranges & Beta-carotene (can interfere with bilirubin measurement) \\
Tryptophan-rich foods & Can cause false-positive urine protein \\
\bottomrule
\end{longtable}

\subsubsection{What This Means For You}

If your doctor orders a fasting test, \textbf{actually fast}. Drinking a latte with milk before a ``fasting'' lipid panel or glucose test will give you a falsely abnormal result, which can lead to unnecessary medications or anxiety. If you accidentally ate before a fasting test, tell your doctor --- the test should be repeated.

\clearpage
\section{Medications That Fool Your Lab Tests}

This is one of the most underappreciated sources of lab error. Thousands of medications affect lab tests, and most patients and many doctors do not think to check.

\subsection{Contraindications and Precautions}

\begin{longtable}{@{}p{4.5cm}p{10.5cm}@{}}
\toprule
\textbf{Medication} & \textbf{Tests Affected} \\
\midrule
Aspirin & Platelet function (prolonged), uric acid (can decrease) \\
Acetaminophen (Tylenol) & Bilirubin (can increase), uric acid (can decrease) \\
Vitamin C (high dose) & Glucose (glucometer falsely low), creatinine (falsely low), occult blood (false negative) \\
Biotin (high dose, $>$5 mg/day) & Troponin, TSH, FT4, FT3, ferritin, PTH (falsely low or high) \\
NSAIDs (ibuprofen, naproxen) & Platelet function, creatinine (can elevate), potassium (can elevate) \\
Oral contraceptives & Clotting factors (elevated), triglycerides (can increase), SHBG (can increase) \\
Corticosteroids (prednisone) & Glucose (elevated), WBC (elevated), cortisol (can suppress), eosinophils (decreased) \\
Antibiotics (fluoroquinolones) & Creatinine (can interfere with assay), QT interval \\
Statins (atorvastatin, etc.) & Liver enzymes (AST, ALT can elevate), CK (can elevate, rarely rhabdomyolysis) \\
Diuretics (furosemide, HCTZ) & Potassium (can decrease), sodium, magnesium, calcium, glucose \\
ACE inhibitors (lisinopril, etc.) & Potassium (can elevate), creatinine (can elevate) \\
Metformin & Vitamin B12 (can decrease long-term) \\
Lithium & Must be monitored (therapeutic range is narrow) \\
Warfarin (Coumadin) & PT/INR (increased --- this is the point), must be monitored \\
Heparin & aPTT (increased --- this is the point) \\
Phenytoin & Liver enzymes, calcium (falsely low on some assays), folate \\
Methyldopa & Direct Coombs test (positive) \\
Levodopa & Uric acid (can increase), liver enzymes \\
\bottomrule
\end{longtable}

\subsubsection{What This Means For You}

\textbf{Always tell your doctor about every medication and supplement you take.} This includes over-the-counter drugs, vitamins, herbal supplements, and recreational substances. Your doctor needs this information to interpret your lab results correctly.

The biotin issue deserves special emphasis. Biotin (vitamin B7) is a common ingredient in hair, skin, and nail supplements. High doses can cause dangerously incorrect results on immunoassays --- your troponin could read falsely low (missing a heart attack), your TSH could read falsely low (faking hyperthyroidism), or your FT4 could read falsely abnormal. \textbf{Stop biotin supplements at least 24--72 hours before any lab work.}

\clearpage
\section{Exercise and Physical Activity}

That morning run before your blood draw is not as harmless as you think.

\begin{longtable}{@{}p{4.5cm}p{10.5cm}@{}}
\toprule
\textbf{Effect of Exercise} & \textbf{Tests Affected} \\
\midrule
Increased potassium & Potassium (falsely elevated) \\
Increased WBC & WBC count (elevated, especially neutrophils) \\
Increased CK & Creatine kinase (falsely elevated, can mimic muscle disease) \\
Increased lactate & Lactate (elevated) \\
Increased AST & AST (elevated, can mimic liver damage) \\
Increased BUN & BUN (elevated from dehydration) \\
Increased glucose & Glucose (elevated from stress hormones) \\
Increased cortisol & Cortisol (elevated) \\
Decreased iron & Serum iron (falsely low) \\
Hemolysis & Potassium (falsely elevated from RBC breakdown), many tests \\
\bottomrule
\end{longtable}

\subsubsection{What This Means For You}

\textbf{Avoid vigorous exercise for at least 24 hours before blood work.} A brisk walk is fine. A hard gym session, a long run, or physically demanding work can elevate your CK, potassium, WBC, and liver enzymes --- all of which can be mistaken for disease.

\clearpage
\section{Posture and How You Sit}

This one surprises most people. Simply sitting versus lying down can change some of your lab values.

\begin{longtable}{@{}p{4.5cm}p{10.5cm}@{}}
\toprule
\textbf{Effect of Upright Posture} & \textbf{Tests Affected} \\
\midrule
Hemoconcentration & Protein, albumin, enzymes (can increase 5--10\%) \\
Decreased plasma volume & Hemoglobin, hematocrit (can increase) \\
Potassium shift & Potassium (can increase from cell stretch) \\
Renin-aldosterone activation & Renin, aldosterone (can increase) \\
\bottomrule
\end{longtable}

\subsubsection{What This Means For You}

If you are getting blood drawn, it is best to sit or lie down for 5--15 minutes before the stick. Standing up with your arm hanging down while waiting for the phlebotomist can shift some values. This is a minor effect for most tests, but it can matter for precise measurements.

\clearpage
\section{Time of Day --- Your Body Has a Clock}

Many hormones and some lab values fluctuate predictably throughout the day.

\begin{longtable}{@{}p{4.5cm}p{10.5cm}@{}}
\toprule
\textbf{Test} & \textbf{Time-of-Day Effect} \\
\midrule
Cortisol & Peaks in early morning (6--8 AM), lowest at midnight \\
ACTH & Peaks in early morning \\
TSH & Peaks in the evening, lowest in the afternoon \\
Iron & Highest in the morning, lowest in the evening \\
PTH & Diurnal variation \\
Testosterone & Highest in the morning \\
FSH, LH & Pulsatile secretion, varies throughout day \\
Growth hormone & Peaks during deep sleep \\
Reticulocyte count & Diurnal variation \\
\bottomrule
\end{longtable}

\subsubsection{What This Means For You}

For consistency, \textbf{have your blood drawn at the same time of day each time} if you are being monitored for a condition. If your cortisol was drawn at 3 PM and it looks low, that may be completely normal --- cortisol is supposed to be low in the afternoon. If your TSH was drawn at 2 PM and it looks borderline, it might be higher if drawn in the evening.

\clearpage
\section{Hemolysis --- When Your Blood Sample Is Damaged}

Hemolysis means your red blood cells have broken open during or after the blood draw, spilling their contents into the serum. This is one of the most common causes of lab error.

\begin{longtable}{@{}p{4.5cm}p{10.5cm}@{}}
\toprule
\textbf{What Hemolysis Does} & \textbf{Why} \\
\midrule
Falsely elevates potassium & RBCs contain 20--40$\times$ more potassium than serum \\
Falsely elevates LDH & LDH is concentrated inside red blood cells \\
Falsely elevates AST & AST is present in red blood cells \\
Falsely elevates magnesium & Magnesium is present in red blood cells \\
Falsely elevates phosphate & Phosphate is present in red blood cells \\
Falsely decreases iron & Hemoglobin interferes with the assay \\
Masks bilirubin & The red color of hemoglobin interferes with measurement \\
\bottomrule
\end{longtable}

\subsubsection{What This Means For You}

If your potassium comes back high and you feel fine, the first thing your doctor should do is check whether the sample was hemolyzed. A hemolyzed sample is \textbf{not reliable} and the test should be redrawn. You have every right to ask: ``Was this sample hemolyzed?'' If it was, insist on a redraw before any treatment is started.

Common causes of hemolysis:
\begin{itemize}
    \item Using a needle that is too small
    \item Drawing blood too quickly (too much suction)
    \item The tourniquet was on too long
    \item The tube was shaken instead of gently inverted
    \item The sample sat too long before processing
    \item Difficult vein access (multiple sticks)
\end{itemize}

\clearpage
\section{Lipemia --- When Your Blood Is Fatty}

Lipemia means your blood serum looks milky or cloudy because of high levels of triglycerides (fats) in your blood.

\begin{longtable}{@{}p{4.5cm}p{10.5cm}@{}}
\toprule
\textbf{Effect of Lipemia} & \textbf{Tests Affected} \\
\midrule
Dilution effect & Falsely decreases the concentration of many analytes \\
Turbidity & Interferes with spectrophotometric assays \\
Falsely increases glucose & Turbidity adds signal to some glucose methods \\
Falsely increases sodium & Some methods are affected by the lipid volume \\
\bottomrule
\end{longtable}

\subsubsection{What This Means For You}

Lipemia is most commonly caused by eating a high-fat meal before the blood draw, but it can also indicate a lipid metabolism disorder. If your sample is lipemic, some tests may need to be performed on the ultracentrifuged sample or redrawn after fasting.

\clearpage
\section{Icterus --- When Your Blood Is Too Yellow}

An icteric sample means your serum is visibly yellow due to very high bilirubin levels. This can interfere with colorimetric assays --- tests that measure color changes.

\subsubsection{What This Means For You}

If you are visibly jaundiced (yellow skin, yellow eyes), your doctor should be aware that some of your lab results may be affected by the high bilirubin. Tests commonly affected include bilirubin itself, amylase, creatinine, and uric acid.

\clearpage
\section{The Tourniquet Problem}

Leaving the tourniquet on for too long (more than 1 minute) during a blood draw can artificially change your results.

\begin{longtable}{@{}p{4.5cm}p{10.5cm}@{}}
\toprule
\textbf{Effect} & \textbf{Tests Affected} \\
\midrule
Hemoconcentration & Proteins, enzymes, cell counts all increase \\
Potassium leakage & Potassium increases from stretch on RBCs \\
Calcium decrease & Calcium decreases as it shifts into cells \\
LDH increase & LDH released from damaged cells \\
\bottomrule
\end{longtable}

\subsubsection{What This Means For You}

If your blood draw takes a long time because the phlebotomist is having trouble finding a vein, or if the tourniquet is on for more than a minute, some values may be affected. This is another reason to speak up if something seems wrong during your blood draw.

\clearpage
\section{Biotin --- The Supplement That Destroys Your Lab Results}

Biotin (vitamin B7) deserves its own section because it is extremely common and its effects are dangerous.

\begin{longtable}{@{}p{4.5cm}p{10.5cm}@{}}
\toprule
\textbf{Effect} & \textbf{Tests Affected} \\
\midrule
Falsely low & TSH, FT4, troponin, PTH, ferritin, NT-proBNP \\
Falsely high & Free T3 (competitive assays) \\
\bottomrule
\end{longtable}

\subsubsection{What This Means For You}

Biotin is found in many supplements marketed for hair, skin, and nail health. It is also a component of many multivitamins. High-dose biotin can cause your troponin to read falsely low (your doctor could miss a heart attack), your TSH to read falsely low (your doctor could diagnose hyperthyroidism you do not have), and your PTH to read falsely low (your doctor could miss parathyroid disease).

\textbf{Stop taking biotin-containing supplements at least 72 hours before any lab work.} Tell your doctor you take biotin. Most labs now add a biotin interference warning to reports when levels are high.

\clearpage
\section{IV Contrast and Imaging Dyes}

If you have had a CT scan or other imaging study with IV contrast dye in the past 48 hours, some of your lab results may be affected.

\begin{longtable}{@{}p{4.5cm}p{10.5cm}@{}}
\toprule
\textbf{Effect} & \textbf{Tests Affected} \\
\midrule
Falsely increases creatinine & Can mimic acute kidney injury \\
Iodide interference & Thyroid function tests \\
\bottomrule
\end{longtable}

\subsubsection{What This Means For You}

If your creatinine is slightly elevated after a CT scan with contrast, it may be the contrast dye, not kidney disease. Your doctor should know about any recent imaging studies when interpreting your results.

\clearpage
\section{How to Get Accurate Lab Results}

Here is your checklist for getting the most accurate lab results possible:

\begin{enumerate}
    \item \textbf{Fast properly.} 8--12 hours for fasting tests. Water only. No coffee with cream.
    \item \textbf{Avoid vigorous exercise for 24 hours} before blood work.
    \item \textbf{Tell your doctor about every medication and supplement} you take, including over-the-counter drugs and vitamins.
    \item \textbf{Stop biotin supplements 72 hours} before lab work.
    \item \textbf{Get blood drawn in the morning} when possible, and at the same time of day for monitoring tests.
    \item \textbf{Sit or lie down for 5--15 minutes} before the blood draw.
    \item \textbf{Do not clench your fist} during the blood draw (this can affect potassium).
    \item \textbf{Ask if your sample was hemolyzed} if you get an unexpected high potassium result.
    \item \textbf{If you have had IV contrast in the past 48 hours, tell your doctor.}
    \item \textbf{If a result does not make sense, ask for the test to be repeated} before starting treatment.
\end{enumerate}

\clearpage
\section{What to Do If You Suspect a False Result}

If your lab result does not match how you feel, or if a result comes back dramatically different from previous results without a clear explanation, here is what to do:

\begin{enumerate}
    \item \textbf{Ask your doctor:} ``Could anything have affected this result? Was the sample hemolyzed?''
    \item \textbf{Request a repeat test} before any treatment is started.
    \item \textbf{Ask which sample tube was used} and whether the correct tube was collected.
    \item \textbf{If the repeat result is the same,} accept that the result is real and work with your doctor to understand why.
    \item \textbf{If the repeat result is different,} the first result was likely the error.
\end{enumerate}

Your lab results are powerful tools, but they are only as good as the sample and the conditions under which they were collected. You are your own best advocate for making sure your results are accurate.
